\documentclass[11pt,letterpaper]{article}
\usepackage[T1]{fontenc}
\usepackage{lmodern,amsmath,amssymb,amsthm,microtype}
\usepackage[margin=1in]{geometry}
\usepackage[hidelinks]{hyperref}
\newtheorem{theorem}{Theorem}
\title{No mixed unit edges in a three-factor triangular lattice sum}
\author{Aafi Khan Malek\\\small Independent Researcher, Ahmedabad, India\\
\small\href{mailto:aafimalek2032@gmail.com}{aafimalek2032@gmail.com}}
\date{November 2026}
\begin{document}
\maketitle
Instead of searching a finite point cloud, we describe the entire
unit-distance graph of one additive coordinate group. This supplies
a rigorous exclusion of a family of proposed constructions, not
a new chromatic bound for the plane.

Let \(s=\sqrt{-3}\), \(r=\sqrt{33}\), and \(h=\sqrt{-7}\),
with positive real or positive imaginary part as appropriate.
Set \(\omega=(1+s)/2\), \(\Lambda=\mathbb Z[\omega]\),
\(\alpha=(r+s)/6\), \(\beta=(3+h)/4\), and
\[
 M=\Lambda+\alpha\Lambda+\beta\Lambda\subset\mathbb C.
\]
\begin{theorem}
The unit-distance graph on \(M\) is the Cartesian product of
three triangular lattice graphs. Its chromatic number is three.
Any chosen distinct nonadjacent pair can be given the same color
in a proper four-coloring of the whole graph.
\end{theorem}
\begin{proof}
The square classes of \(-3,33,-7\) are independent: their
exponents at 11 and 7 detect the latter two classes, and
\(-3\) is not a rational square. Thus \(1,r,h,rh\) are a basis
over \(\mathbb Q(s)\) in this multiquadratic field.
A displacement \(d=t+\alpha u+\beta v\) with \(t,u,v\in\Lambda\)
satisfies
\[
 d=A+rB+hC,\qquad A=t+3v/4+su/6,\quad B=u/6,\quad C=v/4.
\]
Uniqueness of \((t,u,v)\) follows from independence of
\(1,r,h\) over \(\mathbb Q(s)\). Write \(N(z)=z\overline z\).
If \(N(d)=1\), expand:
\[
\begin{split}
 N(d)={}&N(A)+33N(B)+7N(C)
   +r(A\overline B+B\overline A)\\
 &+h(C\overline A-A\overline C)
   +rh(C\overline B-B\overline C).
\end{split}
\]
The last three terms lie in the independent rational directions
\(r,hs,rhs\), respectively, and therefore vanish. In real planar
terms this means that \(A\) is perpendicular to \(B\), and the
oriented areas of the pairs \((C,A),(C,B)\) are zero.

Suppose first that \(B\) and \(C\) are nonzero. Then \(A\)
is both parallel and perpendicular to \(B\), so \(A=0\).
The identity \(12t+9v+2su=0\) gives \(v\in2\Lambda\) upon
reduction modulo \(2\Lambda\). Consequently
\(7N(v)/16\ge7/4>1\), contradicting the norm expansion.

Suppose \(B\ne0\), \(C=0\). The bound \(33N(u)/36\le1\)
forces \(N(u)=1\), because nonzero Eisenstein norms are positive
integers. The condition \(A\perp B\) makes
\(t\overline u\) purely imaginary. Since the purely imaginary
elements of \(\Lambda\) are precisely the integer multiples of
\(s\), write \(t\overline u=ks\) for \(k\in\mathbb Z\).
The remaining norm equation is
\(3(k+1/6)^2=1/12\), whose only integral solution is \(k=0\).
Thus \(t=0\), and the step is entirely in the \(\alpha\) factor.

Suppose \(B=0\), \(C\ne0\). Now \(7N(v)/16\le1\), so
\(N(v)=1\) or 2. The norm \(N(a+b\omega)=a^2+ab+b^2\)
is congruent to \((a-b)^2\) modulo three, ruling out 2.
The parallel condition gives \(t\overline v\in\Lambda\cap\mathbb R
=\mathbb Z\); let this integer be \(k\). The norm equation becomes
\((k+3/4)^2=9/16\), again forcing \(k=0\), so \(t=0\).
This is a step in the \(\beta\) factor alone.

If \(B=C=0\) the step is in \(\Lambda\) alone. Conversely,
every unit step in a single factor is a unit displacement in
\(M\), since \(|\alpha|=|\beta|=1\). This proves the product
description. In each factor color \((a,b)\) by \(a+2b\) modulo
three, and sum the three colors modulo three. Each unit lattice
step changes the color; a triangle in one factor proves that
three colors are required. To give any specified nonadjacent pair
equal colors, recolor both with one fresh fourth color.
\end{proof}

\paragraph{Exact finite example.}
Let \(W\) be 0 and the sixth roots of unity. The proposed
fixture \((W+W)+\alpha W+\beta W\) has 931 distinct exact sites
and 5,250 unit edges. The accompanying standard-library replay
reconstructs all coordinates and edges and checks the
three-coloring and explicit equal-terminal colorings on the
marked nonedge \(0,2\). In particular, this fixture does not
supply a five-color forced-difference gadget.

Enlarging the additive fixture anywhere within \(M\) cannot
supply such a gadget either. A product rotation such as
\(\alpha\beta\) allows a direction outside \(M\); whether such a
different construction helps remains open.
This written note and its verification were developed with AI
assistance. Neither is a Lean proof or a whole-plane result.
\end{document}
